\documentclass[12pt,a4paper]{article}
\usepackage{fontspec}
\usepackage{polyglossia}
\setdefaultlanguage{russian}
\setmainfont[Ligatures=TeX]{Times New Roman}
\newfontfamily\cyrillicfont[Script=Cyrillic,Ligatures=TeX]{Times New Roman}
\newfontfamily\cyrillicfonttt[Script=Cyrillic]{Courier New}
\usepackage{amsmath,amssymb,mathtools}
\usepackage{graphicx,array,multirow,longtable,tabularx}
\usepackage{tikz}
\usetikzlibrary{arrows.meta,calc}
\usepackage{geometry,float,caption,enumitem,indentfirst,titlesec,pdflscape}
\usepackage[hidelinks]{hyperref}
\geometry{left=25mm,right=15mm,top=20mm,bottom=20mm}
\setlength{\parindent}{1.25cm}
\setlength{\parskip}{3pt}
\linespread{1.08}
\setlength{\emergencystretch}{2em}
\captionsetup{labelsep=endash,font=small,justification=centering}
\addto\captionsrussian{\renewcommand{\figurename}{График}}
\titleformat{\section}{\bfseries\normalsize}{\thesection.}{0.5em}{}
\titleformat{\subsection}{\bfseries\normalsize}{\thesubsection.}{0.5em}{}
\titlespacing*{\section}{0pt}{10pt}{5pt}
\titlespacing*{\subsection}{0pt}{7pt}{4pt}
\setlist[enumerate]{nosep,leftmargin=1.2cm}
\newcommand{\formtitle}[2]{\noindent\hfill Форма #1\par\begin{center}\textbf{#2}\end{center}}
\newcommand{\concl}[1]{\par\noindent\textit{Вывод.} #1\par}
\newcommand{\fig}[3][0.94]{\begin{figure}[H]\centering\includegraphics[width=#1\textwidth]{figures/#2}\caption{#3}\end{figure}}
\newcommand{\mm}{\widetilde m}
\newcommand{\DD}{\widetilde D}
\newcommand{\ssig}{\widetilde\sigma}
\hypersetup{pdftitle={Марковские модели систем массового обслуживания},pdfauthor={Якименко Владислав Игоревич}}
\begin{document}
\begin{titlepage}
\centering
Федеральное государственное автономное образовательное учреждение\\
высшего образования\\
«Национальный исследовательский университет ИТМО»\\[5mm]
Факультет программной инженерии и компьютерной техники
\vfill
{\large\bfseries Учебно-исследовательская работа № 2}\\[5mm]
{\large\bfseries «Марковские модели систем массового обслуживания»}\\[5mm]
по дисциплине «Моделирование»\\[3mm]
Вариант 16/16/13
\vfill
\begin{flushright}
\begin{tabular}{ll}
Выполнил: & Якименко Владислав Игоревич\\
Группа: & P3313\\
Проверил: & Авксентьева Елена Юрьевна
\end{tabular}
\end{flushright}
\vfill
Санкт-Петербург\\2026
\end{titlepage}
\setcounter{page}{2}
\section{Постановка задачи и исходные данные}
Цель: разработать марковские модели двух СМО с однородным потоком заявок, рассчитать стационарные характеристики и выбрать вариант с максимальной производительностью.
Принят вариант $N1/N2/NG=16/16/13$. По строке 16 таблицы 1 методички и строке 13 таблицы 2:
\begin{center}\small\renewcommand{\arraystretch}{1.4}
\begin{tabular}{|l|c|c|}\hline
Параметр&СИСТЕМА\_1&СИСТЕМА\_2\\\hline
Число приборов П&3&2\\\hline
ЕН по заданию&$1/1/0$&$2/2$\\\hline
ЕН в расчётной модели&$1/1/0$&$1/2$\\\hline
Закон обслуживания&$M$ во всех П&$H_{2{,}5}$ в П1; $M$ в П2\\\hline
$\lambda$, с$^{-1}$&$0{,}8$&$0{,}8$\\\hline
$b$, с&$5$&$5$\\\hline
Вероятности направления&$0{,}30/0{,}45/0{,}25$&$0{,}30/0{,}70$\\\hline
\end{tabular}\end{center}
Поскольку $N1+N2=32$ --- чётное число, критерий выбирается по $N2$: (а) максимальная производительность. Для двух приборов вероятность направления к П2 равна $0{,}45+0{,}25=0{,}70$.

\section{Описание исследуемых систем}
Входящий поток --- простейший. Заявка направляется к прибору с заданной вероятностью; при заполненном накопителе она теряется. Перенаправления к другому прибору нет. Дисциплина обслуживания --- FIFO.
\begin{gather}
\lambda_j=\lambda\alpha_j,\qquad \mu=\frac1b=0{,}2\;\text{с}^{-1},\\
\boldsymbol\lambda^{(1)}=(0{,}24;0{,}36;0{,}20),\qquad
\boldsymbol\lambda^{(2)}=(0{,}24;0{,}56)\;\text{с}^{-1}.
\end{gather}
Здесь $\alpha_j$ --- вероятность направления к П$j$; $p_s$ далее обозначает вероятность состояния.

В СИСТЕМЕ\_2 время обслуживания в П1 представлено двумя экспоненциальными фазами. Исходная ёмкость $2/2$ даёт $(1+2\cdot3)\cdot4=28$ состояний. По замечанию на с.~7 методички допускается минимальное сокращение ёмкости при числе состояний более 20. Уменьшение первого накопителя с 2 до 1 даёт $(1+2\cdot2)\cdot4=20$ состояний.
\newpage

Для $H_{\nu}$ при $b=5$ с и $\nu=2{,}5$:
\begin{align}
q&=\frac12\left(1-\sqrt{\frac{\nu^2-1}{\nu^2+1}}\right)=0{,}074519,\\
\mu_a&=\frac{2q}{b}=0{,}029807\;\text{с}^{-1},\\
\mu_b&=\frac{2(1-q)}{b}=0{,}370193\;\text{с}^{-1},\\
f(t)&=q\mu_ae^{-\mu_at}+(1-q)\mu_be^{-\mu_bt},\qquad t\geq0.
\end{align}
Фаза $a$ --- медленная, $b$ --- быстрая. Фаза выбирается при начале каждого обслуживания, в том числе после освобождения прибора при непустой очереди:
\begin{gather}
\mathrm M[T]=\frac q{\mu_a}+\frac{1-q}{\mu_b}=5\;\text{с},\\
\mathrm D[T]=\frac{2q}{\mu_a^2}+\frac{2(1-q)}{\mu_b^2}-b^2=156{,}25\;\text{с}^2,\qquad
\frac{\sqrt{\mathrm D[T]}}b=2{,}5.
\end{gather}

\section{Перечень состояний марковского процесса}
\subsection{СИСТЕМА\_1}
Состояние $s=(n_1,n_2,n_3)$ задаёт число заявок у каждого прибора, включая обслуживаемую:
\begin{equation}
\mathcal S_1=\{0,1,2\}\times\{0,1,2\}\times\{0,1\},\qquad |\mathcal S_1|=18.
\end{equation}
При наличии места переход $n_j\to n_j+1$ имеет интенсивность $\lambda_j$; при $n_j>0$ переход $n_j\to n_j-1$ имеет интенсивность $\mu$.

\subsection{СИСТЕМА\_2}
Состояние $s=(s_1,n_2)$ учитывает число заявок и текущую фазу П1:
\begin{equation}
\mathcal S_2=\{0,1_a,1_b,2_a,2_b\}\times\{0,1,2,3\},\qquad |\mathcal S_2|=20.
\end{equation}
Здесь $1_a,1_b$ --- одна заявка на обслуживании; $2_a,2_b$ --- одна на обслуживании и одна в очереди. В таблице формы 1 ниже указан полный перечень состояний в порядке строк матриц.
\newpage

\section{Результаты работы}
\subsection{Размеченные графы переходов}
Каждая вершина $E_i$ соответствует состоянию из формы 1. На дугах указаны интенсивности; потеря заявки не изменяет состояние, поэтому отдельной дугой не изображается.
\begin{figure}[H]\centering
\resizebox{0.96\textwidth}{!}{%
\begin{tikzpicture}[>=Stealth, state/.style={draw,circle,fill=white,inner sep=1pt,minimum size=7mm,font=\footnotesize}, every edge/.style={draw,->,thin}, every node/.style={font=\scriptsize}]
\node[state] (E1) at (0.0,0.0) {$E_{1}$};
\node[state] (E2) at (1.35,-1.15) {$E_{2}$};
\node[state] (E3) at (0.0,-3.5) {$E_{3}$};
\node[state] (E4) at (1.35,-4.65) {$E_{4}$};
\node[state] (E5) at (0.0,-7.0) {$E_{5}$};
\node[state] (E6) at (1.35,-8.15) {$E_{6}$};
\node[state] (E7) at (4.6,0.0) {$E_{7}$};
\node[state] (E8) at (5.949999999999999,-1.15) {$E_{8}$};
\node[state] (E9) at (4.6,-3.5) {$E_{9}$};
\node[state] (E10) at (5.949999999999999,-4.65) {$E_{10}$};
\node[state] (E11) at (4.6,-7.0) {$E_{11}$};
\node[state] (E12) at (5.949999999999999,-8.15) {$E_{12}$};
\node[state] (E13) at (9.2,0.0) {$E_{13}$};
\node[state] (E14) at (10.549999999999999,-1.15) {$E_{14}$};
\node[state] (E15) at (9.2,-3.5) {$E_{15}$};
\node[state] (E16) at (10.549999999999999,-4.65) {$E_{16}$};
\node[state] (E17) at (9.2,-7.0) {$E_{17}$};
\node[state] (E18) at (10.549999999999999,-8.15) {$E_{18}$};
\path (E1) edge[bend left=14] node[midway,sloped,above,fill=white,inner sep=0.6pt] {$\lambda_{3}$} (E2);
\path (E1) edge[bend left=14] node[midway,sloped,above,fill=white,inner sep=0.6pt] {$\lambda_{2}$} (E3);
\path (E1) edge[bend left=14] node[midway,sloped,above,fill=white,inner sep=0.6pt] {$\lambda_{1}$} (E7);
\path (E2) edge[bend left=14] node[midway,sloped,above,fill=white,inner sep=0.6pt] {$\mu$} (E1);
\path (E2) edge[bend left=14] node[midway,sloped,above,fill=white,inner sep=0.6pt] {$\lambda_{2}$} (E4);
\path (E2) edge[bend left=14] node[midway,sloped,above,fill=white,inner sep=0.6pt] {$\lambda_{1}$} (E8);
\path (E3) edge[bend left=14] node[midway,sloped,above,fill=white,inner sep=0.6pt] {$\mu$} (E1);
\path (E3) edge[bend left=14] node[midway,sloped,above,fill=white,inner sep=0.6pt] {$\lambda_{3}$} (E4);
\path (E3) edge[bend left=14] node[midway,sloped,above,fill=white,inner sep=0.6pt] {$\lambda_{2}$} (E5);
\path (E3) edge[bend left=14] node[midway,sloped,above,fill=white,inner sep=0.6pt] {$\lambda_{1}$} (E9);
\path (E4) edge[bend left=14] node[midway,sloped,above,fill=white,inner sep=0.6pt] {$\mu$} (E2);
\path (E4) edge[bend left=14] node[midway,sloped,above,fill=white,inner sep=0.6pt] {$\mu$} (E3);
\path (E4) edge[bend left=14] node[midway,sloped,above,fill=white,inner sep=0.6pt] {$\lambda_{2}$} (E6);
\path (E4) edge[bend left=14] node[midway,sloped,above,fill=white,inner sep=0.6pt] {$\lambda_{1}$} (E10);
\path (E5) edge[bend left=14] node[midway,sloped,above,fill=white,inner sep=0.6pt] {$\mu$} (E3);
\path (E5) edge[bend left=14] node[midway,sloped,above,fill=white,inner sep=0.6pt] {$\lambda_{3}$} (E6);
\path (E5) edge[bend left=14] node[midway,sloped,above,fill=white,inner sep=0.6pt] {$\lambda_{1}$} (E11);
\path (E6) edge[bend left=14] node[midway,sloped,above,fill=white,inner sep=0.6pt] {$\mu$} (E4);
\path (E6) edge[bend left=14] node[midway,sloped,above,fill=white,inner sep=0.6pt] {$\mu$} (E5);
\path (E6) edge[bend left=14] node[midway,sloped,above,fill=white,inner sep=0.6pt] {$\lambda_{1}$} (E12);
\path (E7) edge[bend left=14] node[midway,sloped,above,fill=white,inner sep=0.6pt] {$\mu$} (E1);
\path (E7) edge[bend left=14] node[midway,sloped,above,fill=white,inner sep=0.6pt] {$\lambda_{3}$} (E8);
\path (E7) edge[bend left=14] node[midway,sloped,above,fill=white,inner sep=0.6pt] {$\lambda_{2}$} (E9);
\path (E7) edge[bend left=14] node[midway,sloped,above,fill=white,inner sep=0.6pt] {$\lambda_{1}$} (E13);
\path (E8) edge[bend left=14] node[midway,sloped,above,fill=white,inner sep=0.6pt] {$\mu$} (E2);
\path (E8) edge[bend left=14] node[midway,sloped,above,fill=white,inner sep=0.6pt] {$\mu$} (E7);
\path (E8) edge[bend left=14] node[midway,sloped,above,fill=white,inner sep=0.6pt] {$\lambda_{2}$} (E10);
\path (E8) edge[bend left=14] node[midway,sloped,above,fill=white,inner sep=0.6pt] {$\lambda_{1}$} (E14);
\path (E9) edge[bend left=14] node[midway,sloped,above,fill=white,inner sep=0.6pt] {$\mu$} (E3);
\path (E9) edge[bend left=14] node[midway,sloped,above,fill=white,inner sep=0.6pt] {$\mu$} (E7);
\path (E9) edge[bend left=14] node[midway,sloped,above,fill=white,inner sep=0.6pt] {$\lambda_{3}$} (E10);
\path (E9) edge[bend left=14] node[midway,sloped,above,fill=white,inner sep=0.6pt] {$\lambda_{2}$} (E11);
\path (E9) edge[bend left=14] node[midway,sloped,above,fill=white,inner sep=0.6pt] {$\lambda_{1}$} (E15);
\path (E10) edge[bend left=14] node[midway,sloped,above,fill=white,inner sep=0.6pt] {$\mu$} (E4);
\path (E10) edge[bend left=14] node[midway,sloped,above,fill=white,inner sep=0.6pt] {$\mu$} (E8);
\path (E10) edge[bend left=14] node[midway,sloped,above,fill=white,inner sep=0.6pt] {$\mu$} (E9);
\path (E10) edge[bend left=14] node[midway,sloped,above,fill=white,inner sep=0.6pt] {$\lambda_{2}$} (E12);
\path (E10) edge[bend left=14] node[midway,sloped,above,fill=white,inner sep=0.6pt] {$\lambda_{1}$} (E16);
\path (E11) edge[bend left=14] node[midway,sloped,above,fill=white,inner sep=0.6pt] {$\mu$} (E5);
\path (E11) edge[bend left=14] node[midway,sloped,above,fill=white,inner sep=0.6pt] {$\mu$} (E9);
\path (E11) edge[bend left=14] node[midway,sloped,above,fill=white,inner sep=0.6pt] {$\lambda_{3}$} (E12);
\path (E11) edge[bend left=14] node[midway,sloped,above,fill=white,inner sep=0.6pt] {$\lambda_{1}$} (E17);
\path (E12) edge[bend left=14] node[midway,sloped,above,fill=white,inner sep=0.6pt] {$\mu$} (E6);
\path (E12) edge[bend left=14] node[midway,sloped,above,fill=white,inner sep=0.6pt] {$\mu$} (E10);
\path (E12) edge[bend left=14] node[midway,sloped,above,fill=white,inner sep=0.6pt] {$\mu$} (E11);
\path (E12) edge[bend left=14] node[midway,sloped,above,fill=white,inner sep=0.6pt] {$\lambda_{1}$} (E18);
\path (E13) edge[bend left=14] node[midway,sloped,above,fill=white,inner sep=0.6pt] {$\mu$} (E7);
\path (E13) edge[bend left=14] node[midway,sloped,above,fill=white,inner sep=0.6pt] {$\lambda_{3}$} (E14);
\path (E13) edge[bend left=14] node[midway,sloped,above,fill=white,inner sep=0.6pt] {$\lambda_{2}$} (E15);
\path (E14) edge[bend left=14] node[midway,sloped,above,fill=white,inner sep=0.6pt] {$\mu$} (E8);
\path (E14) edge[bend left=14] node[midway,sloped,above,fill=white,inner sep=0.6pt] {$\mu$} (E13);
\path (E14) edge[bend left=14] node[midway,sloped,above,fill=white,inner sep=0.6pt] {$\lambda_{2}$} (E16);
\path (E15) edge[bend left=14] node[midway,sloped,above,fill=white,inner sep=0.6pt] {$\mu$} (E9);
\path (E15) edge[bend left=14] node[midway,sloped,above,fill=white,inner sep=0.6pt] {$\mu$} (E13);
\path (E15) edge[bend left=14] node[midway,sloped,above,fill=white,inner sep=0.6pt] {$\lambda_{3}$} (E16);
\path (E15) edge[bend left=14] node[midway,sloped,above,fill=white,inner sep=0.6pt] {$\lambda_{2}$} (E17);
\path (E16) edge[bend left=14] node[midway,sloped,above,fill=white,inner sep=0.6pt] {$\mu$} (E10);
\path (E16) edge[bend left=14] node[midway,sloped,above,fill=white,inner sep=0.6pt] {$\mu$} (E14);
\path (E16) edge[bend left=14] node[midway,sloped,above,fill=white,inner sep=0.6pt] {$\mu$} (E15);
\path (E16) edge[bend left=14] node[midway,sloped,above,fill=white,inner sep=0.6pt] {$\lambda_{2}$} (E18);
\path (E17) edge[bend left=14] node[midway,sloped,above,fill=white,inner sep=0.6pt] {$\mu$} (E11);
\path (E17) edge[bend left=14] node[midway,sloped,above,fill=white,inner sep=0.6pt] {$\mu$} (E15);
\path (E17) edge[bend left=14] node[midway,sloped,above,fill=white,inner sep=0.6pt] {$\lambda_{3}$} (E18);
\path (E18) edge[bend left=14] node[midway,sloped,above,fill=white,inner sep=0.6pt] {$\mu$} (E12);
\path (E18) edge[bend left=14] node[midway,sloped,above,fill=white,inner sep=0.6pt] {$\mu$} (E16);
\path (E18) edge[bend left=14] node[midway,sloped,above,fill=white,inner sep=0.6pt] {$\mu$} (E17);
\end{tikzpicture}
}
\caption{Размеченный граф переходов СИСТЕМЫ\_1}
\end{figure}
\begin{equation*}
\lambda_1=0{,}24,\quad \lambda_2=0{,}36,\quad \lambda_3=0{,}20,\quad\mu=0{,}20\;\text{с}^{-1}.
\end{equation*}
\concl{Полный граф содержит 18 состояний. Число заявок меняется на единицу только у одного прибора за переход.}
\clearpage
\begin{landscape}
\begin{figure}[H]\centering
\resizebox{0.98\linewidth}{!}{%
\begin{tikzpicture}[>=Stealth, state/.style={draw,circle,fill=white,inner sep=1pt,minimum size=7mm,font=\footnotesize}, every edge/.style={draw,->,thin}, every node/.style={font=\scriptsize}]
\node[state] (E1) at (0.0,0) {$E_{1}$};
\node[state] (E2) at (5.8,0) {$E_{2}$};
\node[state] (E3) at (11.6,0) {$E_{3}$};
\node[state] (E4) at (17.4,0) {$E_{4}$};
\node[state] (E5) at (-1.2,-2.3) {$E_{5}$};
\node[state] (E6) at (4.6,-2.3) {$E_{6}$};
\node[state] (E7) at (10.4,-2.3) {$E_{7}$};
\node[state] (E8) at (16.2,-2.3) {$E_{8}$};
\node[state] (E9) at (1.2,-3.8) {$E_{9}$};
\node[state] (E10) at (7.0,-3.8) {$E_{10}$};
\node[state] (E11) at (12.799999999999999,-3.8) {$E_{11}$};
\node[state] (E12) at (18.599999999999998,-3.8) {$E_{12}$};
\node[state] (E13) at (-1.2,-6.2) {$E_{13}$};
\node[state] (E14) at (4.6,-6.2) {$E_{14}$};
\node[state] (E15) at (10.4,-6.2) {$E_{15}$};
\node[state] (E16) at (16.2,-6.2) {$E_{16}$};
\node[state] (E17) at (1.2,-7.7) {$E_{17}$};
\node[state] (E18) at (7.0,-7.7) {$E_{18}$};
\node[state] (E19) at (12.799999999999999,-7.7) {$E_{19}$};
\node[state] (E20) at (18.599999999999998,-7.7) {$E_{20}$};
\path (E1) edge[bend left=14] node[midway,sloped,above,fill=white,inner sep=0.6pt] {$\lambda_2$} (E2);
\path (E1) edge[bend left=14] node[midway,sloped,above,fill=white,inner sep=0.6pt] {$\lambda_1q$} (E5);
\path (E1) edge[bend left=14] node[midway,sloped,above,fill=white,inner sep=0.6pt] {$\lambda_1(1-q)$} (E9);
\path (E2) edge[bend left=14] node[midway,sloped,above,fill=white,inner sep=0.6pt] {$\mu$} (E1);
\path (E2) edge[bend left=14] node[midway,sloped,above,fill=white,inner sep=0.6pt] {$\lambda_2$} (E3);
\path (E2) edge[bend left=14] node[midway,sloped,above,fill=white,inner sep=0.6pt] {$\lambda_1q$} (E6);
\path (E2) edge[bend left=14] node[midway,sloped,above,fill=white,inner sep=0.6pt] {$\lambda_1(1-q)$} (E10);
\path (E3) edge[bend left=14] node[midway,sloped,above,fill=white,inner sep=0.6pt] {$\mu$} (E2);
\path (E3) edge[bend left=14] node[midway,sloped,above,fill=white,inner sep=0.6pt] {$\lambda_2$} (E4);
\path (E3) edge[bend left=14] node[midway,sloped,above,fill=white,inner sep=0.6pt] {$\lambda_1q$} (E7);
\path (E3) edge[bend left=14] node[midway,sloped,above,fill=white,inner sep=0.6pt] {$\lambda_1(1-q)$} (E11);
\path (E4) edge[bend left=14] node[midway,sloped,above,fill=white,inner sep=0.6pt] {$\mu$} (E3);
\path (E4) edge[bend left=14] node[midway,sloped,above,fill=white,inner sep=0.6pt] {$\lambda_1q$} (E8);
\path (E4) edge[bend left=14] node[midway,sloped,above,fill=white,inner sep=0.6pt] {$\lambda_1(1-q)$} (E12);
\path (E5) edge[bend left=14] node[midway,sloped,above,fill=white,inner sep=0.6pt] {$\mu_a$} (E1);
\path (E5) edge[bend left=14] node[midway,sloped,above,fill=white,inner sep=0.6pt] {$\lambda_2$} (E6);
\path (E5) edge[bend left=14] node[midway,sloped,above,fill=white,inner sep=0.6pt] {$\lambda_1$} (E13);
\path (E6) edge[bend left=14] node[midway,sloped,above,fill=white,inner sep=0.6pt] {$\mu_a$} (E2);
\path (E6) edge[bend left=14] node[midway,sloped,above,fill=white,inner sep=0.6pt] {$\mu$} (E5);
\path (E6) edge[bend left=14] node[midway,sloped,above,fill=white,inner sep=0.6pt] {$\lambda_2$} (E7);
\path (E6) edge[bend left=14] node[midway,sloped,above,fill=white,inner sep=0.6pt] {$\lambda_1$} (E14);
\path (E7) edge[bend left=14] node[midway,sloped,above,fill=white,inner sep=0.6pt] {$\mu_a$} (E3);
\path (E7) edge[bend left=14] node[midway,sloped,above,fill=white,inner sep=0.6pt] {$\mu$} (E6);
\path (E7) edge[bend left=14] node[midway,sloped,above,fill=white,inner sep=0.6pt] {$\lambda_2$} (E8);
\path (E7) edge[bend left=14] node[midway,sloped,above,fill=white,inner sep=0.6pt] {$\lambda_1$} (E15);
\path (E8) edge[bend left=14] node[midway,sloped,above,fill=white,inner sep=0.6pt] {$\mu_a$} (E4);
\path (E8) edge[bend left=14] node[midway,sloped,above,fill=white,inner sep=0.6pt] {$\mu$} (E7);
\path (E8) edge[bend left=14] node[midway,sloped,above,fill=white,inner sep=0.6pt] {$\lambda_1$} (E16);
\path (E9) edge[bend left=14] node[midway,sloped,above,fill=white,inner sep=0.6pt] {$\mu_b$} (E1);
\path (E9) edge[bend left=14] node[midway,sloped,above,fill=white,inner sep=0.6pt] {$\lambda_2$} (E10);
\path (E9) edge[bend left=14] node[midway,sloped,above,fill=white,inner sep=0.6pt] {$\lambda_1$} (E17);
\path (E10) edge[bend left=14] node[midway,sloped,above,fill=white,inner sep=0.6pt] {$\mu_b$} (E2);
\path (E10) edge[bend left=14] node[midway,sloped,above,fill=white,inner sep=0.6pt] {$\mu$} (E9);
\path (E10) edge[bend left=14] node[midway,sloped,above,fill=white,inner sep=0.6pt] {$\lambda_2$} (E11);
\path (E10) edge[bend left=14] node[midway,sloped,above,fill=white,inner sep=0.6pt] {$\lambda_1$} (E18);
\path (E11) edge[bend left=14] node[midway,sloped,above,fill=white,inner sep=0.6pt] {$\mu_b$} (E3);
\path (E11) edge[bend left=14] node[midway,sloped,above,fill=white,inner sep=0.6pt] {$\mu$} (E10);
\path (E11) edge[bend left=14] node[midway,sloped,above,fill=white,inner sep=0.6pt] {$\lambda_2$} (E12);
\path (E11) edge[bend left=14] node[midway,sloped,above,fill=white,inner sep=0.6pt] {$\lambda_1$} (E19);
\path (E12) edge[bend left=14] node[midway,sloped,above,fill=white,inner sep=0.6pt] {$\mu_b$} (E4);
\path (E12) edge[bend left=14] node[midway,sloped,above,fill=white,inner sep=0.6pt] {$\mu$} (E11);
\path (E12) edge[bend left=14] node[midway,sloped,above,fill=white,inner sep=0.6pt] {$\lambda_1$} (E20);
\path (E13) edge[bend left=14] node[midway,sloped,above,fill=white,inner sep=0.6pt] {$\mu_aq$} (E5);
\path (E13) edge[bend left=23] node[midway,sloped,above,fill=white,inner sep=0.6pt] {$\mu_a(1-q)$} (E9);
\path (E13) edge[bend left=14] node[midway,sloped,above,fill=white,inner sep=0.6pt] {$\lambda_2$} (E14);
\path (E14) edge[bend left=14] node[midway,sloped,above,fill=white,inner sep=0.6pt] {$\mu_aq$} (E6);
\path (E14) edge[bend left=23] node[midway,sloped,above,fill=white,inner sep=0.6pt] {$\mu_a(1-q)$} (E10);
\path (E14) edge[bend left=14] node[midway,sloped,above,fill=white,inner sep=0.6pt] {$\mu$} (E13);
\path (E14) edge[bend left=14] node[midway,sloped,above,fill=white,inner sep=0.6pt] {$\lambda_2$} (E15);
\path (E15) edge[bend left=14] node[midway,sloped,above,fill=white,inner sep=0.6pt] {$\mu_aq$} (E7);
\path (E15) edge[bend left=23] node[midway,sloped,above,fill=white,inner sep=0.6pt] {$\mu_a(1-q)$} (E11);
\path (E15) edge[bend left=14] node[midway,sloped,above,fill=white,inner sep=0.6pt] {$\mu$} (E14);
\path (E15) edge[bend left=14] node[midway,sloped,above,fill=white,inner sep=0.6pt] {$\lambda_2$} (E16);
\path (E16) edge[bend left=14] node[midway,sloped,above,fill=white,inner sep=0.6pt] {$\mu_aq$} (E8);
\path (E16) edge[bend left=23] node[midway,sloped,above,fill=white,inner sep=0.6pt] {$\mu_a(1-q)$} (E12);
\path (E16) edge[bend left=14] node[midway,sloped,above,fill=white,inner sep=0.6pt] {$\mu$} (E15);
\path (E17) edge[bend left=23] node[midway,sloped,above,fill=white,inner sep=0.6pt] {$\mu_bq$} (E5);
\path (E17) edge[bend left=14] node[midway,sloped,above,fill=white,inner sep=0.6pt] {$\mu_b(1-q)$} (E9);
\path (E17) edge[bend left=14] node[midway,sloped,above,fill=white,inner sep=0.6pt] {$\lambda_2$} (E18);
\path (E18) edge[bend left=23] node[midway,sloped,above,fill=white,inner sep=0.6pt] {$\mu_bq$} (E6);
\path (E18) edge[bend left=14] node[midway,sloped,above,fill=white,inner sep=0.6pt] {$\mu_b(1-q)$} (E10);
\path (E18) edge[bend left=14] node[midway,sloped,above,fill=white,inner sep=0.6pt] {$\mu$} (E17);
\path (E18) edge[bend left=14] node[midway,sloped,above,fill=white,inner sep=0.6pt] {$\lambda_2$} (E19);
\path (E19) edge[bend left=23] node[midway,sloped,above,fill=white,inner sep=0.6pt] {$\mu_bq$} (E7);
\path (E19) edge[bend left=14] node[midway,sloped,above,fill=white,inner sep=0.6pt] {$\mu_b(1-q)$} (E11);
\path (E19) edge[bend left=14] node[midway,sloped,above,fill=white,inner sep=0.6pt] {$\mu$} (E18);
\path (E19) edge[bend left=14] node[midway,sloped,above,fill=white,inner sep=0.6pt] {$\lambda_2$} (E20);
\path (E20) edge[bend left=23] node[midway,sloped,above,fill=white,inner sep=0.6pt] {$\mu_bq$} (E8);
\path (E20) edge[bend left=14] node[midway,sloped,above,fill=white,inner sep=0.6pt] {$\mu_b(1-q)$} (E12);
\path (E20) edge[bend left=14] node[midway,sloped,above,fill=white,inner sep=0.6pt] {$\mu$} (E19);
\end{tikzpicture}
}
\caption{Размеченный граф переходов СИСТЕМЫ\_2}
\end{figure}
\begin{gather*}
\lambda_1=0{,}24,\quad\lambda_2=0{,}56,\quad\mu=0{,}20,\quad
q=0{,}074519,\quad \mu_a=0{,}029807,\quad\mu_b=0{,}370193.\\
1_a\xrightarrow{\mu_a}0,\quad1_b\xrightarrow{\mu_b}0,\quad
2_{\varphi}\xrightarrow{\mu_{\varphi}q}1_a,\quad
2_{\varphi}\xrightarrow{\mu_{\varphi}(1-q)}1_b,\quad\varphi\in\{a,b\}.
\end{gather*}
\concl{Фазовая кодировка делает длительность пребывания в каждом состоянии экспоненциальной. Получена цепь Маркова с 20 состояниями.}
\end{landscape}
\clearpage

\subsection{Матрицы интенсивностей переходов}
Обозначения соответствуют теории к УИР 2: $G=[g_{ij}]$ --- матрица интенсивностей. Для $i\ne j$ значение $g_{ij}$ равно интенсивности перехода $E_i\to E_j$:
\begin{equation}
g_{ii}=-\sum_{\substack{j=1\\j\ne i}}^{n}g_{ij},\qquad
\sum_{j=1}^{n}g_{ij}=0.\label{eq:G}
\end{equation}
В расчётной книге эти же матрицы обозначены $Q_1,Q_2$; здесь $G_1=Q_1$, $G_2=Q_2$. Каждый блок ниже содержит все строки и половину столбцов. Числа округлены только при выводе; решение получено по неокруглённым матрицам.

\subsection{Система уравнений}
Стационарные вероятности найдены из системы
\begin{equation}
\begin{cases}
\displaystyle\sum_{i=1}^{n}p_i g_{ij}=0,& j=1,\ldots,n,\\[2mm]
\displaystyle\sum_{i=1}^{n}p_i=1,&p_i\geq0.
\end{cases}\label{eq:balance}
\end{equation}
При численном решении одно зависимое балансовое уравнение заменено нормировкой.

Методичка предусматривает расчёт в программе MARK. Приведённые ниже числа получены независимым решением системы~\eqref{eq:balance} в Python; выполнение расчёта в MARK не подтверждено.
\newpage
\subsection{Матрица $G_1$}
\noindent Столбцы $E_{1},\ldots,E_{9}$; значения в $\text{с}^{-1}$.
\begin{center}\scriptsize\setlength{\tabcolsep}{2.6pt}\renewcommand{\arraystretch}{1.18}
\begin{tabular}{|c|r|r|r|r|r|r|r|r|r|}\hline
Из / в&$E_{1}$&$E_{2}$&$E_{3}$&$E_{4}$&$E_{5}$&$E_{6}$&$E_{7}$&$E_{8}$&$E_{9}$\\\hline
$E_{1}$&$-0{,}80000$&$0{,}20000$&$0{,}36000$&$0$&$0$&$0$&$0{,}24000$&$0$&$0$\\\hline
$E_{2}$&$0{,}20000$&$-0{,}80000$&$0$&$0{,}36000$&$0$&$0$&$0$&$0{,}24000$&$0$\\\hline
$E_{3}$&$0{,}20000$&$0$&$-1{,}00000$&$0{,}20000$&$0{,}36000$&$0$&$0$&$0$&$0{,}24000$\\\hline
$E_{4}$&$0$&$0{,}20000$&$0{,}20000$&$-1{,}00000$&$0$&$0{,}36000$&$0$&$0$&$0$\\\hline
$E_{5}$&$0$&$0$&$0{,}20000$&$0$&$-0{,}64000$&$0{,}20000$&$0$&$0$&$0$\\\hline
$E_{6}$&$0$&$0$&$0$&$0{,}20000$&$0{,}20000$&$-0{,}64000$&$0$&$0$&$0$\\\hline
$E_{7}$&$0{,}20000$&$0$&$0$&$0$&$0$&$0$&$-1{,}00000$&$0{,}20000$&$0{,}36000$\\\hline
$E_{8}$&$0$&$0{,}20000$&$0$&$0$&$0$&$0$&$0{,}20000$&$-1{,}00000$&$0$\\\hline
$E_{9}$&$0$&$0$&$0{,}20000$&$0$&$0$&$0$&$0{,}20000$&$0$&$-1{,}20000$\\\hline
$E_{10}$&$0$&$0$&$0$&$0{,}20000$&$0$&$0$&$0$&$0{,}20000$&$0{,}20000$\\\hline
$E_{11}$&$0$&$0$&$0$&$0$&$0{,}20000$&$0$&$0$&$0$&$0{,}20000$\\\hline
$E_{12}$&$0$&$0$&$0$&$0$&$0$&$0{,}20000$&$0$&$0$&$0$\\\hline
$E_{13}$&$0$&$0$&$0$&$0$&$0$&$0$&$0{,}20000$&$0$&$0$\\\hline
$E_{14}$&$0$&$0$&$0$&$0$&$0$&$0$&$0$&$0{,}20000$&$0$\\\hline
$E_{15}$&$0$&$0$&$0$&$0$&$0$&$0$&$0$&$0$&$0{,}20000$\\\hline
$E_{16}$&$0$&$0$&$0$&$0$&$0$&$0$&$0$&$0$&$0$\\\hline
$E_{17}$&$0$&$0$&$0$&$0$&$0$&$0$&$0$&$0$&$0$\\\hline
$E_{18}$&$0$&$0$&$0$&$0$&$0$&$0$&$0$&$0$&$0$\\\hline
\end{tabular}\end{center}
\noindent Столбцы $E_{10},\ldots,E_{18}$; значения в $\text{с}^{-1}$.
\begin{center}\scriptsize\setlength{\tabcolsep}{2.6pt}\renewcommand{\arraystretch}{1.18}
\begin{tabular}{|c|r|r|r|r|r|r|r|r|r|}\hline
Из / в&$E_{10}$&$E_{11}$&$E_{12}$&$E_{13}$&$E_{14}$&$E_{15}$&$E_{16}$&$E_{17}$&$E_{18}$\\\hline
$E_{1}$&$0$&$0$&$0$&$0$&$0$&$0$&$0$&$0$&$0$\\\hline
$E_{2}$&$0$&$0$&$0$&$0$&$0$&$0$&$0$&$0$&$0$\\\hline
$E_{3}$&$0$&$0$&$0$&$0$&$0$&$0$&$0$&$0$&$0$\\\hline
$E_{4}$&$0{,}24000$&$0$&$0$&$0$&$0$&$0$&$0$&$0$&$0$\\\hline
$E_{5}$&$0$&$0{,}24000$&$0$&$0$&$0$&$0$&$0$&$0$&$0$\\\hline
$E_{6}$&$0$&$0$&$0{,}24000$&$0$&$0$&$0$&$0$&$0$&$0$\\\hline
$E_{7}$&$0$&$0$&$0$&$0{,}24000$&$0$&$0$&$0$&$0$&$0$\\\hline
$E_{8}$&$0{,}36000$&$0$&$0$&$0$&$0{,}24000$&$0$&$0$&$0$&$0$\\\hline
$E_{9}$&$0{,}20000$&$0{,}36000$&$0$&$0$&$0$&$0{,}24000$&$0$&$0$&$0$\\\hline
$E_{10}$&$-1{,}20000$&$0$&$0{,}36000$&$0$&$0$&$0$&$0{,}24000$&$0$&$0$\\\hline
$E_{11}$&$0$&$-0{,}84000$&$0{,}20000$&$0$&$0$&$0$&$0$&$0{,}24000$&$0$\\\hline
$E_{12}$&$0{,}20000$&$0{,}20000$&$-0{,}84000$&$0$&$0$&$0$&$0$&$0$&$0{,}24000$\\\hline
$E_{13}$&$0$&$0$&$0$&$-0{,}76000$&$0{,}20000$&$0{,}36000$&$0$&$0$&$0$\\\hline
$E_{14}$&$0$&$0$&$0$&$0{,}20000$&$-0{,}76000$&$0$&$0{,}36000$&$0$&$0$\\\hline
$E_{15}$&$0$&$0$&$0$&$0{,}20000$&$0$&$-0{,}96000$&$0{,}20000$&$0{,}36000$&$0$\\\hline
$E_{16}$&$0{,}20000$&$0$&$0$&$0$&$0{,}20000$&$0{,}20000$&$-0{,}96000$&$0$&$0{,}36000$\\\hline
$E_{17}$&$0$&$0{,}20000$&$0$&$0$&$0$&$0{,}20000$&$0$&$-0{,}60000$&$0{,}20000$\\\hline
$E_{18}$&$0$&$0$&$0{,}20000$&$0$&$0$&$0$&$0{,}20000$&$0{,}20000$&$-0{,}60000$\\\hline
\end{tabular}\end{center}
\newpage
\subsection{Матрица $G_2$}
\noindent Столбцы $E_{1},\ldots,E_{10}$; значения в $\text{с}^{-1}$.
\begin{center}\scriptsize\setlength{\tabcolsep}{2.6pt}\renewcommand{\arraystretch}{1.18}
\begin{tabular}{|c|r|r|r|r|r|r|r|r|r|r|}\hline
Из / в&$E_{1}$&$E_{2}$&$E_{3}$&$E_{4}$&$E_{5}$&$E_{6}$&$E_{7}$&$E_{8}$&$E_{9}$&$E_{10}$\\\hline
$E_{1}$&$-0{,}80000$&$0{,}56000$&$0$&$0$&$0{,}01788$&$0$&$0$&$0$&$0{,}22212$&$0$\\\hline
$E_{2}$&$0{,}20000$&$-1{,}00000$&$0{,}56000$&$0$&$0$&$0{,}01788$&$0$&$0$&$0$&$0{,}22212$\\\hline
$E_{3}$&$0$&$0{,}20000$&$-1{,}00000$&$0{,}56000$&$0$&$0$&$0{,}01788$&$0$&$0$&$0$\\\hline
$E_{4}$&$0$&$0$&$0{,}20000$&$-0{,}44000$&$0$&$0$&$0$&$0{,}01788$&$0$&$0$\\\hline
$E_{5}$&$0{,}02981$&$0$&$0$&$0$&$-0{,}82981$&$0{,}56000$&$0$&$0$&$0$&$0$\\\hline
$E_{6}$&$0$&$0{,}02981$&$0$&$0$&$0{,}20000$&$-1{,}02981$&$0{,}56000$&$0$&$0$&$0$\\\hline
$E_{7}$&$0$&$0$&$0{,}02981$&$0$&$0$&$0{,}20000$&$-1{,}02981$&$0{,}56000$&$0$&$0$\\\hline
$E_{8}$&$0$&$0$&$0$&$0{,}02981$&$0$&$0$&$0{,}20000$&$-0{,}46981$&$0$&$0$\\\hline
$E_{9}$&$0{,}37019$&$0$&$0$&$0$&$0$&$0$&$0$&$0$&$-1{,}17019$&$0{,}56000$\\\hline
$E_{10}$&$0$&$0{,}37019$&$0$&$0$&$0$&$0$&$0$&$0$&$0{,}20000$&$-1{,}37019$\\\hline
$E_{11}$&$0$&$0$&$0{,}37019$&$0$&$0$&$0$&$0$&$0$&$0$&$0{,}20000$\\\hline
$E_{12}$&$0$&$0$&$0$&$0{,}37019$&$0$&$0$&$0$&$0$&$0$&$0$\\\hline
$E_{13}$&$0$&$0$&$0$&$0$&$0{,}00222$&$0$&$0$&$0$&$0{,}02759$&$0$\\\hline
$E_{14}$&$0$&$0$&$0$&$0$&$0$&$0{,}00222$&$0$&$0$&$0$&$0{,}02759$\\\hline
$E_{15}$&$0$&$0$&$0$&$0$&$0$&$0$&$0{,}00222$&$0$&$0$&$0$\\\hline
$E_{16}$&$0$&$0$&$0$&$0$&$0$&$0$&$0$&$0{,}00222$&$0$&$0$\\\hline
$E_{17}$&$0$&$0$&$0$&$0$&$0{,}02759$&$0$&$0$&$0$&$0{,}34261$&$0$\\\hline
$E_{18}$&$0$&$0$&$0$&$0$&$0$&$0{,}02759$&$0$&$0$&$0$&$0{,}34261$\\\hline
$E_{19}$&$0$&$0$&$0$&$0$&$0$&$0$&$0{,}02759$&$0$&$0$&$0$\\\hline
$E_{20}$&$0$&$0$&$0$&$0$&$0$&$0$&$0$&$0{,}02759$&$0$&$0$\\\hline
\end{tabular}\end{center}
\noindent Столбцы $E_{11},\ldots,E_{20}$; значения в $\text{с}^{-1}$.
\begin{center}\scriptsize\setlength{\tabcolsep}{2.6pt}\renewcommand{\arraystretch}{1.18}
\begin{tabular}{|c|r|r|r|r|r|r|r|r|r|r|}\hline
Из / в&$E_{11}$&$E_{12}$&$E_{13}$&$E_{14}$&$E_{15}$&$E_{16}$&$E_{17}$&$E_{18}$&$E_{19}$&$E_{20}$\\\hline
$E_{1}$&$0$&$0$&$0$&$0$&$0$&$0$&$0$&$0$&$0$&$0$\\\hline
$E_{2}$&$0$&$0$&$0$&$0$&$0$&$0$&$0$&$0$&$0$&$0$\\\hline
$E_{3}$&$0{,}22212$&$0$&$0$&$0$&$0$&$0$&$0$&$0$&$0$&$0$\\\hline
$E_{4}$&$0$&$0{,}22212$&$0$&$0$&$0$&$0$&$0$&$0$&$0$&$0$\\\hline
$E_{5}$&$0$&$0$&$0{,}24000$&$0$&$0$&$0$&$0$&$0$&$0$&$0$\\\hline
$E_{6}$&$0$&$0$&$0$&$0{,}24000$&$0$&$0$&$0$&$0$&$0$&$0$\\\hline
$E_{7}$&$0$&$0$&$0$&$0$&$0{,}24000$&$0$&$0$&$0$&$0$&$0$\\\hline
$E_{8}$&$0$&$0$&$0$&$0$&$0$&$0{,}24000$&$0$&$0$&$0$&$0$\\\hline
$E_{9}$&$0$&$0$&$0$&$0$&$0$&$0$&$0{,}24000$&$0$&$0$&$0$\\\hline
$E_{10}$&$0{,}56000$&$0$&$0$&$0$&$0$&$0$&$0$&$0{,}24000$&$0$&$0$\\\hline
$E_{11}$&$-1{,}37019$&$0{,}56000$&$0$&$0$&$0$&$0$&$0$&$0$&$0{,}24000$&$0$\\\hline
$E_{12}$&$0{,}20000$&$-0{,}81019$&$0$&$0$&$0$&$0$&$0$&$0$&$0$&$0{,}24000$\\\hline
$E_{13}$&$0$&$0$&$-0{,}58981$&$0{,}56000$&$0$&$0$&$0$&$0$&$0$&$0$\\\hline
$E_{14}$&$0$&$0$&$0{,}20000$&$-0{,}78981$&$0{,}56000$&$0$&$0$&$0$&$0$&$0$\\\hline
$E_{15}$&$0{,}02759$&$0$&$0$&$0{,}20000$&$-0{,}78981$&$0{,}56000$&$0$&$0$&$0$&$0$\\\hline
$E_{16}$&$0$&$0{,}02759$&$0$&$0$&$0{,}20000$&$-0{,}22981$&$0$&$0$&$0$&$0$\\\hline
$E_{17}$&$0$&$0$&$0$&$0$&$0$&$0$&$-0{,}93019$&$0{,}56000$&$0$&$0$\\\hline
$E_{18}$&$0$&$0$&$0$&$0$&$0$&$0$&$0{,}20000$&$-1{,}13019$&$0{,}56000$&$0$\\\hline
$E_{19}$&$0{,}34261$&$0$&$0$&$0$&$0$&$0$&$0$&$0{,}20000$&$-1{,}13019$&$0{,}56000$\\\hline
$E_{20}$&$0$&$0{,}34261$&$0$&$0$&$0$&$0$&$0$&$0$&$0{,}20000$&$-0{,}57019$\\\hline
\end{tabular}\end{center}

\newpage
\subsection{Стационарные вероятности состояний}
\formtitle{1}{Стационарные вероятности состояний}
\begin{center}\small\setlength{\tabcolsep}{9pt}\renewcommand{\arraystretch}{1.15}
\begin{tabular}{|c|c|r|c|r|}\hline
\multirow{2}{*}{\shortstack{Номер\\состояния}}&\multicolumn{2}{c|}{СИСТЕМА\_1}&\multicolumn{2}{c|}{СИСТЕМА\_2}\\\cline{2-5}
&Обозн.&Вер-ть&Обозн.&Вер-ть\\\hline
1&$(0,0,0)$&$0{,}022742$&$(0,0)$&$0{,}009583$\\\hline
2&$(0,0,1)$&$0{,}022742$&$(0,1)$&$0{,}026833$\\\hline
3&$(0,1,0)$&$0{,}040936$&$(0,2)$&$0{,}075133$\\\hline
4&$(0,1,1)$&$0{,}040936$&$(0,3)$&$0{,}210372$\\\hline
5&$(0,2,0)$&$0{,}073685$&$(1_{a},0)$&$0{,}001115$\\\hline
6&$(0,2,1)$&$0{,}073685$&$(1_{a},1)$&$0{,}003122$\\\hline
7&$(1,0,0)$&$0{,}027291$&$(1_{a},2)$&$0{,}008742$\\\hline
8&$(1,0,1)$&$0{,}027291$&$(1_{a},3)$&$0{,}024477$\\\hline
9&$(1,1,0)$&$0{,}049123$&$(1_{b},0)$&$0{,}006123$\\\hline
10&$(1,1,1)$&$0{,}049123$&$(1_{b},1)$&$0{,}017145$\\\hline
11&$(1,2,0)$&$0{,}088422$&$(1_{b},2)$&$0{,}048006$\\\hline
12&$(1,2,1)$&$0{,}088422$&$(1_{b},3)$&$0{,}134416$\\\hline
13&$(2,0,0)$&$0{,}032749$&$(2_{a},0)$&$0{,}008978$\\\hline
14&$(2,0,1)$&$0{,}032749$&$(2_{a},1)$&$0{,}025138$\\\hline
15&$(2,1,0)$&$0{,}058948$&$(2_{a},2)$&$0{,}070386$\\\hline
16&$(2,1,1)$&$0{,}058948$&$(2_{a},3)$&$0{,}197082$\\\hline
17&$(2,2,0)$&$0{,}106106$&$(2_{b},0)$&$0{,}003970$\\\hline
18&$(2,2,1)$&$0{,}106106$&$(2_{b},1)$&$0{,}011115$\\\hline
19&---&---&$(2_{b},2)$&$0{,}031123$\\\hline
20&---&---&$(2_{b},3)$&$0{,}087143$\\\hline
\end{tabular}\end{center}

\begin{equation*}
\sum_i p_i^{(1)}=1,\qquad\sum_i p_i^{(2)}=1,\qquad
\max_j\left|\sum_i p_i g_{ij}\right|<10^{-14}.
\end{equation*}
\concl{Решения нормированы и удовлетворяют балансовым уравнениям. Максимумы: $p_{17}^{(1)}=p_{18}^{(1)}=0{,}106106$; $p_4^{(2)}=0{,}210372$.}
\newpage

\subsection{Формулы расчёта характеристик}
Обозначим $K_j=r_j+1$ --- вместимость прибора с накопителем, $n_j(s)$ --- число заявок у П$j$ в состоянии $s$, $A_j=\{s:n_j(s)<K_j\}$ --- принимающие состояния, $(x)^+=\max(x,0)$.
В таблице формы 2 все характеристики вычисляются через $p_s$; формулы Литтла используются только при проверке.

Для поступающей заявки условное среднее ожидание $d_j(s)$ в принимающем состоянии равно
\begin{equation}
d_j(s)=\begin{cases}
0,&n_j(s)=0,\\
n_j(s)b,&n_j(s)>0,\ \text{экспоненциальный П},\\
\dfrac1{\mu_{\varphi(s)}}+(n_j(s)-1)b,&n_j(s)>0,\ H_{\nu}\text{-П}.
\end{cases}\label{eq:waitstate}
\end{equation}
Для П1 СИСТЕМЫ\_2 принимаются только $0,1_a,1_b$, поэтому
\begin{equation}
w_1^{(2)}=\frac{\displaystyle\sum_{s:s_1=1_a}\frac{p_s}{\mu_a}+
\displaystyle\sum_{s:s_1=1_b}\frac{p_s}{\mu_b}}{1-\pi_1^{(2)}}.
\end{equation}
Для остальных приборов:
\begin{equation}
w_j=\frac{\displaystyle\sum_{s\in A_j}n_j(s)b\,p_s}{1-\pi_j},\qquad u_j=w_j+b.
\end{equation}
Усреднение системных времён проводится по принятому потоку:
\begin{equation}
W=\frac{\sum_j\lambda_j^{\prime}w_j}{\Lambda^{\prime}},\qquad
U=\frac{\sum_j\lambda_j^{\prime}u_j}{\Lambda^{\prime}},\qquad
\Lambda^{\prime}=\sum_j\lambda_j^{\prime}.
\end{equation}
При трёх экспоненциальных приборах максимум возможного $R$ равен 3, при двух --- 2. Суммарная загрузка $R=\sum_j\rho_j$ --- среднее число занятых приборов.
\newpage
\formtitle{2}{Характеристики СИСТЕМ}
\begin{center}\small\setlength{\tabcolsep}{3pt}\renewcommand{\arraystretch}{2.0}
\begin{tabular}{|p{27mm}|c|>{\centering\arraybackslash}p{64mm}|r|r|}\hline
Хар-ка&Прибор&Расчётная формула&СИСТ.1&СИСТ.2\\\hline
\multirow{4}{27mm}{\shortstack[l]{Нагрузка}}&П1&$y_j=\lambda_jb$&$1{,}200000$&$1{,}200000$\\\cline{2-5}
&П2&$y_j=\lambda_jb$&$1{,}800000$&$2{,}800000$\\\cline{2-5}
&П3&$y_j=\lambda_jb$&$1{,}000000$&---\\\cline{2-5}
&Сумм.&$Y=\sum_jy_j$&$4{,}000000$&$4{,}000000$\\\hline
\multirow{4}{27mm}{\shortstack[l]{Загрузка}}&П1&$\rho_j=\sum_{s:n_j(s)>0}p_s$&$0{,}725275$&$0{,}678079$\\\cline{2-5}
&П2&$\rho_j=\sum_{s:n_j(s)>0}p_s$&$0{,}834437$&$0{,}970231$\\\cline{2-5}
&П3&$\rho_j=\sum_{s:n_j(s)>0}p_s$&$0{,}500000$&---\\\cline{2-5}
&Сумм.&$R=\sum_j\rho_j$&$2{,}059712$&$1{,}648310$\\\hline
\multirow{4}{27mm}{\shortstack[l]{Длина\\очереди}}&П1&$l_j=\sum_s(n_j(s)-1)^+p_s$&$0{,}395604$&$0{,}434934$\\\cline{2-5}
&П2&$l_j=\sum_s(n_j(s)-1)^+p_s$&$0{,}536424$&$1{,}540367$\\\cline{2-5}
&П3&$l_j=\sum_s(n_j(s)-1)^+p_s$&$0{,}000000$&---\\\cline{2-5}
&Сумм.&$L=\sum_jl_j$&$0{,}932028$&$1{,}975301$\\\hline
\multirow{4}{27mm}{\shortstack[l]{Число\\заявок}}&П1&$m_j=\sum_sn_j(s)p_s$&$1{,}120879$&$1{,}113013$\\\cline{2-5}
&П2&$m_j=\sum_sn_j(s)p_s$&$1{,}370861$&$2{,}510598$\\\cline{2-5}
&П3&$m_j=\sum_sn_j(s)p_s$&$0{,}500000$&---\\\cline{2-5}
&Сумм.&$M=\sum_jm_j$&$2{,}991740$&$3{,}623611$\\\hline
\end{tabular}\end{center}
\noindent\small «---» означает отсутствие П3 в СИСТЕМЕ\_2.\normalsize
\newpage
\formtitle{2}{Характеристики СИСТЕМ (продолжение)}
\begin{center}\small\setlength{\tabcolsep}{3pt}\renewcommand{\arraystretch}{2.0}
\begin{tabular}{|p{27mm}|c|>{\centering\arraybackslash}p{64mm}|r|r|}\hline
Хар-ка&Прибор&Расчётная формула&СИСТ.1&СИСТ.2\\\hline
\multirow{4}{27mm}{\shortstack[l]{Время\\ожидания, с}}&П1&$w_j=\dfrac{\sum_{s\in A_j}p_s d_j(s)}{1-\pi_j}$&$2{,}727273$&$3{,}207105$\\\cline{2-5}
&П2&$w_j=\dfrac{\sum_{s\in A_j}p_s d_j(s)}{1-\pi_j}$&$3{,}214286$&$7{,}938144$\\\cline{2-5}
&П3&$w_j=\dfrac{\sum_{s\in A_j}p_s d_j(s)}{1-\pi_j}$&$0{,}000000$&---\\\cline{2-5}
&Сумм.&$W=\dfrac{\sum_j\lambda_j^{\prime}w_j}{\Lambda^{\prime}}$&$2{,}262521$&$5{,}991897$\\\hline
\multirow{4}{27mm}{\shortstack[l]{Время\\пребывания, с}}&П1&$u_j=w_j+b$&$7{,}727273$&$8{,}207105$\\\cline{2-5}
&П2&$u_j=w_j+b$&$8{,}214286$&$12{,}938144$\\\cline{2-5}
&П3&$u_j=w_j+b$&$5{,}000000$&---\\\cline{2-5}
&Сумм.&$U=\dfrac{\sum_j\lambda_j^{\prime}u_j}{\Lambda^{\prime}}$&$7{,}262521$&$10{,}991897$\\\hline
\multirow{4}{27mm}{\shortstack[l]{Вероятность\\потери}}&П1&$\pi_j=\sum_{s:n_j(s)=K_j}p_s$&$0{,}395604$&$0{,}434934$\\\cline{2-5}
&П2&$\pi_j=\sum_{s:n_j(s)=K_j}p_s$&$0{,}536424$&$0{,}653489$\\\cline{2-5}
&П3&$\pi_j=\sum_{s:n_j(s)=K_j}p_s$&$0{,}500000$&---\\\cline{2-5}
&Сумм.&$\pi=\dfrac{\sum_j\lambda_j\pi_j}{\lambda}$&$0{,}485072$&$0{,}587922$\\\hline
\multirow{4}{27mm}{\shortstack[l]{Производи-\\тельность, с$^{-1}$}}&П1&$\lambda_j^{\prime}=\lambda_j(1-\pi_j)$&$0{,}145055$&$0{,}135616$\\\cline{2-5}
&П2&$\lambda_j^{\prime}=\lambda_j(1-\pi_j)$&$0{,}166887$&$0{,}194046$\\\cline{2-5}
&П3&$\lambda_j^{\prime}=\lambda_j(1-\pi_j)$&$0{,}100000$&---\\\cline{2-5}
&Сумм.&$\Lambda^{\prime}=\sum_j\lambda_j^{\prime}$&$0{,}411942$&$0{,}329662$\\\hline
\end{tabular}\end{center}
\noindent\small «---» означает отсутствие П3 в СИСТЕМЕ\_2.\normalsize

\concl{При $\lambda=0{,}8$ СИСТЕМА\_1 имеет большую производительность и меньшие потери, очередь и среднее время пребывания.}
\newpage

\subsection{Проверка полученных результатов}
Формулы Литтла используются после прямого расчёта:
\begin{equation}
l_j=\lambda_j^{\prime}w_j,\qquad m_j=\lambda_j^{\prime}u_j,\qquad
L=\Lambda^{\prime}W,\qquad M=\Lambda^{\prime}U.
\end{equation}
\begin{center}\small\renewcommand{\arraystretch}{1.35}
\begin{tabular}{|l|r|r|}\hline
Проверка&СИСТЕМА\_1&СИСТЕМА\_2\\\hline
$\sum_i p_i$&$1{,}000000000$&$1{,}000000000$\\\hline
$L$&$0{,}932028237$&$1{,}975300855$\\\hline
$\Lambda^{\prime}W$&$0{,}932028237$&$1{,}975300855$\\\hline
$M$&$2{,}991740048$&$3{,}623610941$\\\hline
$\Lambda^{\prime}U$&$2{,}991740048$&$3{,}623610941$\\\hline
\end{tabular}\end{center}
\concl{Прямой расчёт характеристик согласуется с нормировкой, балансом потоков и формулами Литтла.}

\subsection{Сравнительный анализ}
При изменении $\lambda$ от $0{,}1$ до $2{,}0$ с$^{-1}$ с шагом $0{,}1$ остальные параметры фиксированы. Для каждого значения заново решена система~\eqref{eq:balance}.
\fig{throughput_sensitivity.pdf}{Зависимость производительности от интенсивности потока}
\concl{На исследованной сетке при $\lambda\geq0{,}2$ СИСТЕМА\_1 производительнее. При $\lambda=0{,}1$ небольшое преимущество имеет СИСТЕМА\_2.}
\newpage
\fig{loss_sensitivity.pdf}{Зависимость вероятности потери от интенсивности потока}
\concl{При $\lambda=0{,}8$ потери равны $0{,}485072$ и $0{,}587922$. В СИСТЕМЕ\_2 к П2 направляется $70\%$ потока; его накопитель часто заполнен.}
\fig{sojourn_sensitivity.pdf}{Зависимость среднего времени пребывания от интенсивности потока}
\concl{На исследованной сетке среднее время пребывания в СИСТЕМЕ\_1 меньше. При заданной нагрузке разность равна $10{,}991897-7{,}262521=3{,}729376$ с.}

\section{Выбор наилучшего варианта}
Критерий (а): $\Lambda^{\prime}\to\max$. При $\lambda=0{,}8$ с$^{-1}$:
\begin{gather}
\Lambda_1^{\prime}=0{,}411942\;\text{с}^{-1}>\Lambda_2^{\prime}=0{,}329662\;\text{с}^{-1},\\
\frac{\Lambda_1^{\prime}-\Lambda_2^{\prime}}{\Lambda_2^{\prime}}\cdot100\%=24{,}959\%.
\end{gather}
Выбирается СИСТЕМА\_1. Сравнение относится к расчётной СИСТЕМЕ\_2 с ЕН $1/2$, полученной допустимым уменьшением накопителя; результаты исходной структуры $2/2$ здесь не рассчитывались.
\end{document}
